\section{\bsq{enum} Value Names}\label{sec:EnumValues}
When the \ccjpl had been written in 1999, enums did not yet exist in the programming language. But because \lstinline|enum| values are seen as special cases of constants\index{constant}, the usual assumption is that their names have to follow the same rules as for regular constants\index{field!constant}: the name should be all uppercase with words separated by underscores~(\bdq{\texttt{\_}}).

Basically, this works fine, but when using static imports, name clashes can occur for common terms that can be used in different contexts. So it makes sense to add a prefix to the name of an \lstinline|enum| value\index{enum}, similar to what we did with the categories or groups for regular constants\index{field!constant} (see \hyperref[sec:Constants]{above}). This may look like this:
\keeplisting{12}
\begin{lstlisting}
enum Direction
{
    DIR_LEFT,
    DIR_RIGHT;
}   //  enum Direction

enum Alignment
{
    ALIGN_CENTER,
    ALIGN_LEFT,
    ALIGN_RIGHT;
}   //  enum Alignment
\end{lstlisting}
where \lstinline|DIR_| and \lstinline|ALIGN_| as the respective prefixes, differentiating the \texttt{LEFT} for the direction from that for the alignment.

Using the prefix would allow now to use the values with static imports of their classes; otherwise, they have to be used with their class names as prefix.

\keeplisting{12}
Alternatively, you can use this:
\begin{lstlisting}
enum Direction
{
    DIR_Left,
    DIR_Right;
}   //  enum Direction

enum Alignment
{
    ALIGN_Center,
    ALIGN_Left,
    ALIGN_Right;
}   //  enum Alignment
\end{lstlisting}
Again the \textit{CamelCased} suffix is the \bdq{name} itself, while the \bsq{proper} typed prefix is just the category. This is the recommended approach if the suffix consists of multiple words that otherwise require additional underscores.

\subsubsection{Principles}
\begin{enumerate}[label=P\arabic*.]
	\item{The names of \lstinline|enum| values\index{enum} are all uppercase, with underscores (\bdq{\texttt{\_}}) as separators between words.}
	
	\item{Alternatively, \lstinline|enum| values\index{enum} can be named using a regular name (all uppercase~…) followed by an underscore as prefix, and the name in CamelCase.}
\end{enumerate}

%==============================================================================
% End of file!!
%==============================================================================
